% p182_1; source PDF page 182; manually verified editable table.
\begin{tikzpicture}[x=1cm,y=-1cm,inner sep=0pt,outer sep=0pt]
\path[use as bounding box] (0,0) rectangle (16.000,7.250);
\draw[line width=.4pt] (0.000,0.000) rectangle (2.000,0.800);
\node[align=center,text width=1.740cm,font={\linespread{1}\fontsize{10}{13.50}\selectfont\bfseries}] at (1.000,0.400) {};
\draw[line width=.4pt] (2.000,0.000) rectangle (7.400,0.800);
\node[align=center,text width=5.140cm,font={\linespread{1}\fontsize{10}{13.50}\selectfont\bfseries}] at (4.700,0.400) {社区矫正};
\draw[line width=.4pt] (7.400,0.000) rectangle (16.000,0.800);
\node[align=center,text width=8.340cm,font={\linespread{1}\fontsize{10}{13.50}\selectfont\bfseries}] at (11.700,0.400) {禁止令};
\draw[line width=.4pt] (0.000,0.800) rectangle (2.000,2.150);
\node[align=left,text width=1.740cm,font={\linespread{1}\fontsize{10}{13.50}\selectfont}] at (1.000,1.475) {性质};
\draw[line width=.4pt] (2.000,0.800) rectangle (7.400,2.150);
\node[align=left,text width=5.140cm,font={\linespread{1}\fontsize{10}{13.50}\selectfont}] at (4.700,1.475) {是一种刑罚执行方式，不是新的刑种};
\draw[line width=.4pt] (7.400,0.800) rectangle (16.000,2.150);
\node[align=left,text width=8.340cm,font={\linespread{1}\fontsize{10}{13.50}\selectfont}] at (11.700,1.475) {是一种非刑罚处罚措施，不是刑罚执行方式，不是新的刑种};
\draw[line width=.4pt] (0.000,2.150) rectangle (2.000,3.000);
\node[align=left,text width=1.740cm,font={\linespread{1}\fontsize{10}{13.50}\selectfont}] at (1.000,2.575) {特点};
\draw[line width=.4pt] (2.000,2.150) rectangle (7.400,3.000);
\node[align=left,text width=5.140cm,font={\linespread{1}\fontsize{10}{13.50}\selectfont}] at (4.700,2.575) {非监禁，也即不关押};
\draw[line width=.4pt] (7.400,2.150) rectangle (16.000,3.000);
\node[align=left,text width=8.340cm,font={\linespread{1}\fontsize{10}{13.50}\selectfont}] at (11.700,2.575) {禁止从事某些活动};
\draw[line width=.4pt] (0.000,3.000) rectangle (2.000,4.700);
\node[align=left,text width=1.740cm,font={\linespread{1}\fontsize{10}{13.50}\selectfont}] at (1.000,3.850) {适用领域};
\draw[line width=.4pt] (2.000,3.000) rectangle (7.400,4.700);
\node[align=left,text width=5.140cm,font={\linespread{1}\fontsize{10}{13.50}\selectfont}] at (4.700,3.850) {管制、缓刑、假释，因为这三者均未关押};
\draw[line width=.4pt] (7.400,3.000) rectangle (16.000,4.700);
\node[align=left,text width=8.340cm,font={\linespread{1}\fontsize{10}{13.50}\selectfont}] at (11.700,3.850) {管制、缓刑，不包括假释，因为禁止令是在审判时作出的，而非在执行阶段作出的};
\draw[line width=.4pt] (0.000,4.700) rectangle (2.000,6.400);
\node[align=left,text width=1.740cm,font={\linespread{1}\fontsize{10}{13.50}\selectfont}] at (1.000,5.550) {期限};
\draw[line width=.4pt] (2.000,4.700) rectangle (7.400,6.400);
\node[align=left,text width=5.140cm,font={\linespread{1}\fontsize{10}{13.50}\selectfont}] at (4.700,5.550) {与管制的期限、缓刑考验期的期限、假释考验期的期限相同};
\draw[line width=.4pt] (7.400,4.700) rectangle (16.000,6.400);
\node[align=left,text width=8.340cm,font={\linespread{1}\fontsize{10}{13.50}\selectfont}] at (11.700,5.550) {既可与管制的期限、缓刑考验期的期限相同，也可短于二者};
\draw[line width=.4pt] (0.000,6.400) rectangle (2.000,7.250);
\node[align=left,text width=1.740cm,font={\linespread{1}\fontsize{10}{13.50}\selectfont}] at (1.000,6.825) {执行机构};
\draw[line width=.4pt] (2.000,6.400) rectangle (7.400,7.250);
\node[align=left,text width=5.140cm,font={\linespread{1}\fontsize{10}{13.50}\selectfont}] at (4.700,6.825) {社区矫正机构};
\draw[line width=.4pt] (7.400,6.400) rectangle (16.000,7.250);
\node[align=left,text width=8.340cm,font={\linespread{1}\fontsize{10}{13.50}\selectfont}] at (11.700,6.825) {社区矫正机构};
\end{tikzpicture}
